\documentclass{article}
\usepackage{amsmath}
\usepackage{amssymb}
\usepackage{graphicx}
\usepackage{enumitem}
\usepackage[utf8]{inputenc}
\graphicspath{{Imagenes\}}	
\title{\Huge Taller de Herramientas Computacionales}
\author{Mariana Yasmin Martínez García}
\date{15 de enero del 2019}
\begin{document}
	\maketitle
	\begin{figure}
		\includegraphics[scale=.5]{Imagenes/1}
		\caption{Escudo de la Facultad de ciencias}
		\label{fig:1}
	\end{figure}

\newpage
%\section{Expresiones Matemáticas}
\section*{Expresiones Matemáticas}
$ \alpha + \beta $ \\  % En lugar de los símbolos de peso podemos usar:  \(\)
\[ \alpha + \beta \] %Este lo centra

\section*{Índices y subíndices}
$X_{2}$ \\
$X^{2}$  \\

\section*{Fracciones}
$\frac{3}{5}$\\ \\
$\frac{\frac{3}{4}}{\frac{2}{3}}$ \\ \\

$\sqrt{2} + \sqrt{3^2}^5$ \\ \\
\% \\ \\%para que no se tome como comentario
$\int_{a}^{b} x^2 dx$ \\ \\
$\int_{a}^{b} x^2 \partial x^2$ \\ \\
$3 \quad 2$


\end{document}
